\chapter{Metodología}
La metodología elegida tiene un impacto fundamental en el desarrollo del proyecto,
existen diferentes y variadas opciones, personalmente no fue hasta que comencé a 
indagar en este tema cuando me comencé a sentir seguro de mi trabajo, realizar 
cualquier proyecto era más cuestión de suerte y ganas que de puesta en 
practica del conocimiento. De las muchas opciones, me decanto por dos:
\begin{description}
	\item[RUP] Rational Unified Process.
	\item[XP] Extreme Programming.
\end{description}
\section{Dos es mejor que una}
¿Se puede desarrollar un proyecto con dos metodologías?, la respuesta depende del
tipo de proyecto, en este caso la respuesta es un contundente si, el proyecto es
muy amplio implica partes mecánicas, eléctricas y de programación, mi propuesta 
es utilizar \emph{RUP} para las partes mecánicas y eléctricas y \emph{XP} para 
la programación.
\subsection{RUP}
Dentro de las metodologías hay dos polos, las ágiles y las guiadas por plan, RUP
es una guiada por plan, sin embargo no está en un extremo opuesto a las ágiles.

RUP consiste en una serie de pasos bien definidos, esta diseñada para abarcar 
proyectos muy grandes, proyectos que abarcan colaboraciones entre países o entre
empresas transnacionales, mi proyecto es pequeño en comparación, sin embargo 
considero adecuado aprovechar muchas de las características de esta metodología,
los siguientes escalones son el esqueleto del proceso de desarrollo:
\begin{description}
	\item[Análisis] En este primer paso estamos en el peor de los mundos, pues
		tenemos poca información del problema y una especificación pobre.
		Hacer una análisis en la metodología RUP, implica entender el 
		lenguaje del dominio, no es lo mismo hablar con un médico que con
		un agricultor o un chef de cocina, entender a quién tiene el 
		problema en es un paso fundamental para poder diseñar
		una solución que satisfaga sus necesidades, no entender el 
		lenguaje implicaría una serie de adivinanzas desafortunadas.

		Cuando el entendimiento del lenguaje es suficiente, la prioridad
		es encontrar los requerimientos, nuevamente los requerimientos 
		vienen de parte de quien requiere la solución.

		Finalmente transformar los requerimientos obtenidos en lenguaje 
		de la ingeniería que corresponda da paso a la siguiente etapa.
	\item[Diseño] Con los requerimientos definidos se tiene lo necesario para
		comenzar a diseñar la funcionalidad requerida, en el caso de 
\end{description}<++>
\subsection{XP}
